% Cover page and document-information page.
\begin{titlepage}
  \begin{tikzpicture}[remember picture,overlay]
    \fill[reportnavy] (current page.north west) rectangle
      ([yshift=-17mm]current page.north east);
    \node[anchor=west,text=white,font=\sffamily\small]
      at ([xshift=26mm,yshift=-8.5mm]current page.north west)
      {\addfontfeatures{LetterSpace=10}SCIENTIFIC PROJECT REPORT};
    \fill[reportnavy] (current page.south west) rectangle
      ([yshift=5mm]current page.south east);
  \end{tikzpicture}%
  \vspace*{34mm}

  \begin{flushleft}
    {\bfseries\color{reportnavy}\fontsize{26}{32}\selectfont
      From \textit{Moringa oleifera} Radiobiology\\to Explainable Machine Learning\par}
    \vspace{8mm}
    {\color{reportnavy}\rule{0.32\linewidth}{1.2pt}}\par
    \vspace{8mm}
    {\fontsize{15.5}{21}\selectfont\color{reportgrey}
      Cellular Radiosensitivity Prediction and\\
      Evidence-Based Compound Prioritisation\par}
  \end{flushleft}

  \vfill

  \begin{minipage}[b]{0.82\linewidth}
    \small\color{reportgrey}
    A reproducible investigation linking an experimental study of
    \textit{Moringa oleifera} seed and leaf extracts in irradiated cells to an explainable
    benchmark of clonogenic radiosensitivity across 237 independent cancer cell-line
    components, and to a transparent, evidence-based shortlist of plant constituents for
    future experimental validation.
  \end{minipage}

  \vspace{12mm}
  {\color{reportrule}\rule{\linewidth}{0.5pt}}\par\vspace{2mm}
  {\sffamily\small\color{reportgrey}22 October 2026\par}
\end{titlepage}

% ------------------------------------------------------------ document information
\thispagestyle{empty}
\vspace*{\fill}
\begingroup\small
\noindent{\sffamily\bfseries\color{reportnavy}Document information}\par\smallskip
\noindent\begin{tabular}{@{}L{0.27\linewidth}L{0.69\linewidth}@{}}
  \toprule
  Title & From \textit{Moringa oleifera} Radiobiology to Explainable Machine Learning:
    Cellular Radiosensitivity Prediction and Evidence-Based Compound Prioritisation \\
  Document type & Scientific project report, final edition \\
  Date & 22 October 2026 \\
  Primary analysis & Run \path{primary_v1}: five prespecified models, ten repeats of
    grouped five-fold nested cross-validation, 250 of 250 model-fold tasks completed \\
  Further analyses & Runs \path{exploratory_m3_v1}, \path{loto_v1} and \path{scales_v1};
    explainability and triage of \path{primary_v1}; compound evidence and ranking
    \path{compounds_v1} \\
  Not executed & Negative controls (1,000 draws of each null) and the coverage
    simulation for an interval method (Chapter~\ref{ch:implementation}) \\
  Outcome data licence & CLARIFi v1.0, CC BY-NC 4.0; data DOI
    \doilink{10.5281/zenodo.20710748} and the licence must be cited in any reuse \\
  \bottomrule
\end{tabular}

\medskip
\noindent Every number in this report is taken from a saved output of the project
pipeline; the source file of each result table is named in its caption, and the
correspondence with the project's table identifiers is given in
Table~\ref{tab:provenance}. The figures are the vector graphics produced by the
project's figure scripts. No model was refitted for this report.
\par\endgroup
\clearpage
